\section{The Ratio Difference}

Ratio is not another reporting layer on top of the same kind of engine. It
replaces the engine. The difference reduces to one idea: \textbf{the mechanical
core of accounting is a conservation law, and conservation laws can be proved.}
From that single commitment, five buyer-visible advantages follow.

% Black box (legacy) vs glass box (Ratio): the central marketing contrast.
\begin{figure}[h]
\centering
\resizebox{0.96\textwidth}{!}{%
\begin{tikzpicture}[
  >={Stealth[length=2.4mm]}, font=\sffamily,
  io/.style={font=\sffamily\footnotesize, text=ratioBrown, align=center},
  comp/.style={rounded corners=2pt, draw=ratioTan, fill=ratioCream,
               font=\sffamily\footnotesize, text=ratioInk, align=center,
               minimum height=7mm, text width=30mm, inner sep=2pt},
]
% ── Left: legacy black box ──
\node[rectangle, rounded corners=5pt, fill=ratioInk, minimum width=46mm,
      minimum height=34mm] (lb) at (0,0) {};
\node[font=\sffamily\bfseries, text=white] at (0,1.35) {Incumbent platform};
\node[font=\fontsize{30}{30}\selectfont\bfseries, text=ratioSand] at (0,-0.2) {?};
\node[io, text=ratioInk] at (0,-2.55) {``Trust the output.''};
\node[io] (lin) at (-4.6,0.4) {transactions,\\ prices};
\node[io] (lout) at (4.6,0.4) {reports,\\ NAV, tax lots};
\draw[->, ratioBrown] (lin) -- (lb.west);
\draw[->, ratioBrown] (lb.east) -- (lout);

% ── Right: Ratio glass box ──
\begin{scope}[xshift=11.2cm]
\node[rectangle, rounded corners=5pt, draw=ratioBrown, line width=1pt,
      fill=ratioSand!35, minimum width=50mm, minimum height=34mm] (rb) at (0,0) {};
\node[font=\sffamily\bfseries, text=ratioInk] at (0,1.35) {Ratio};
\node[comp] (k) at (0,0.45) {Conservation kernel \yes\,proven};
\node[comp] (p) at (0,-0.55) {Lean~4 proof \,$\to$\, Rust};
\node[comp] (l) at (0,-1.45) {append-only, integer-exact ledger};
\node[io, text=posGreen] at (0,-2.55) {``Verify the output.''};
\node[io] (rin) at (-5.0,0.4) {transactions,\\ prices};
\node[io] (rout) at (5.0,0.4) {reports +\\ checkable proof};
\draw[->, ratioBrown] (rin) -- (rb.west);
\draw[->, ratioBrown] (rb.east) -- (rout);
\end{scope}
\end{tikzpicture}%
}
\caption*{\small\color{ratioBrown}\sffamily Legacy engines are opaque: identical
inputs, an unseeable middle, and output you take on faith. Ratio's engine is
transparent and proven --- and emits a checkable result, not just a number.}
\end{figure}


\subsection{1. Provable correctness}
Ratio models every transaction as a movement of value that must net to zero
across every dimension --- currencies, positions, tax lots, P\&L. That invariant
is \emph{machine-checked in Lean~4}, and the Rust engine you run is generated
from the proven source. Correctness is established for all inputs by a proof
assistant, not estimated by a test suite. When a ledger is assembled from valid
transactions, it is guaranteed balanced --- as a theorem.

\subsection{2. Reproducible, audit-ready history}
The ledger is an append-only log of integer-exact entries, folded
deterministically. There is no floating-point drift, no order-dependence, no
hidden mutable state. Any past state replays bit-for-bit, under the exact rules
in force then, because those rules are pinned as content-addressed artifacts.
Audit, restatement, and ``as-of'' reporting become deterministic queries.

\subsection{3. AI-native, with the model fenced out of the truth}
Ratio uses LLMs where they are genuinely useful --- turning an advisor's intent
into a precise specification, drafting client-ready explanations --- and nowhere
they are dangerous. Every model output is compiled and checked against the
kernel's invariants before it can touch a ledger. The model proposes; the proven
kernel disposes.

\subsection{4. A core you can audit, and data you can take}
The accounting kernel is authored in Lean~4 and emitted to readable Rust behind
a standard gRPC API, published under a source-available license. Your reviewers,
your auditors, and your regulators can read the arithmetic your firm depends on
and check the theorem for themselves --- which is the guarantee that matters,
and the one no incumbent offers. Your data leaves in open formats, complete,
whenever you want it. Ratio is not permissively licensed and is not open source:
operating the platform is a commercial arrangement with us.

\subsection{5. Performance that is fast \emph{because} it is simple}
A non-expressive kernel over fixed-width integer records is branch-minimal: its
cost is set by memory bandwidth, not by how many business cases the engine must
dispatch on. Exactness removes the reconciliation passes legacy engines spend
cycles on. Simple, proven, and quick are the same design decision.
